%% Generated from AUDIT.md by scripts/md_to_latex.py - edit the Markdown, then re-run the script.
%% Compile with XeLaTeX or LuaLaTeX (e.g. tectonic AUDIT.tex).
\documentclass[10pt,a4paper]{article}
\usepackage[margin=2cm]{geometry}
\usepackage{fontspec}
\setmainfont{DejaVuSerif.ttf}[BoldFont=DejaVuSerif-Bold.ttf, ItalicFont=DejaVuSerif-Italic.ttf,
  BoldItalicFont=DejaVuSerif-BoldItalic.ttf, Scale=0.92]
\setmonofont{DejaVuSansMono.ttf}[BoldFont=DejaVuSansMono-Bold.ttf, Scale=0.85]
\usepackage{graphicx,longtable,array,enumitem,fvextra,xcolor}
\usepackage[hidelinks]{hyperref}
\usepackage{xurl}
\setlist{nosep,leftmargin=1.4em}
\setlength{\parindent}{0pt}
\setlength{\parskip}{0.5em}
\renewcommand{\arraystretch}{1.2}
\begin{document}

{\LARGE\bfseries Cabin\_Noise code audit (2026-11-02)\par}\bigskip

\textbf{Scope:} every module of \texttt{cabin\_\allowbreak{}noise/\allowbreak{}core}, the case runners, the GUI and the docs. The question was whether the code is a faithful low-fidelity interior-noise estimate built from the referenced papers (METHOD.md), with nothing beyond them.

\textbf{Method:}

\begin{itemize}
\item \texttt{pyflakes} over the code and tests;
\item a grep for leftovers of the retired pipeline;
\item a line-by-line read of each block against its paper;
\item tests for every fix;
\item a re-run of all cases and the GUI smoke test.
\end{itemize}

\section*{Verdict}

\textbf{The pipeline is the papers' method.}

{\small
\begin{longtable}{>{\raggedright\arraybackslash}p{0.383\linewidth}>{\raggedright\arraybackslash}p{0.557\linewidth}}
\hline
\textbf{Block} & \textbf{Status} \\
\hline
\endhead
Gutin & matches RP-1258 ch. 1 \\
Goody & matches Goody 2004 \\
CR-132602 eqs. (46) and (52) & match the report, with eq. (53)'s misprint resolved by eq. (47) \\
TM-85722 eq. (1), IL decomposition & match \\
Sabine, Schroeder statistic, IEC weighting & match \\
\hline
\end{longtable}}

Each block has a test against an independent evaluation; see VERIFICATION.md.

The audit found \textbf{no error in any paper equation}. It found and fixed two wiring issues and some housekeeping.

\section*{Findings and decisions}

{\small
\begin{longtable}{>{\raggedright\arraybackslash}p{0.032\linewidth}>{\raggedright\arraybackslash}p{0.371\linewidth}>{\raggedright\arraybackslash}p{0.106\linewidth}>{\raggedright\arraybackslash}p{0.371\linewidth}}
\hline
\textbf{\#} & \textbf{Finding} & \textbf{Severity} & \textbf{Decision} \\
\hline
\endhead
F1 & Window and door area was applied once, whatever the exposed area. The boundary-layer path, which exposes both sides, therefore had half the window share of the tone path. & wiring error & \textbf{Fixed.} Openings are defined per side and scale with the exposure S / side\_area. Test: \texttt{test\_\allowbreak{}openings\_\allowbreak{}scale\_\allowbreak{}with\_\allowbreak{}the\_\allowbreak{}exposed\_\allowbreak{}area}. The end-to-end hand check caught the change and was updated to the documented rule. \\
F2 & The V6 runner runs each engine separately (the pipeline holds one rpm per run) and added the runs' total levels, so the boundary layer was counted twice. & runner error & \textbf{Fixed.} Tones are added across engines and the boundary layer is counted once. Test: \texttt{test\_\allowbreak{}v6\_\allowbreak{}runner\_\allowbreak{}counts\_\allowbreak{}the\_\allowbreak{}boundary\_\allowbreak{}layer\_\allowbreak{}once}. Re-run: V6 is unchanged (372/519, not met), because the doubling appeared in both bare and treated runs and cancelled in the IL. \\
F3 & \texttt{log10(0)} runtime warnings for a static segment with no boundary layer. & cosmetic & Fixed in \texttt{energetic\_\allowbreak{}sum} (an empty spectrum sums to -inf). \\
F4 & Dead code and stale labels: \texttt{log\_\allowbreak{}subfrequencies}, an unused variable, unused GUI imports, the sweep label "first panel", and a V4 quantity naming a runner that does not exist. & housekeeping & Removed or relabelled. \\
F5 & ρc convention: tau uses cabin rho c although the exterior is at altitude. & checked, correct & No change. For a mass-controlled wall the transmitted power depends on the blocked pressure and the receiving side's rho c only, so cabin rho c is right. This also reproduces TM-85722 App. A's 1-2 dB cabin-pressure effect. \\
F6 & Gutin is a static formula but is used with the flight thrust and torque. & method limit & No change. It is the paper's method; documented in METHOD.md and the Propeller tab. \\
F7 & The 95th-percentile seat (+4.8 dB) is applied in all bands, but the statistic holds only above the Schroeder frequency. & method limit & No change to the method. The Interior plot legend now says where it holds. \\
F8 & Both walls' tau use the same rho c, as CR-132602 assumes one medium. & checked & Consistent with F5. \\
\hline
\end{longtable}}

\section*{Not changed (by design)}

\begin{itemize}
\item The validation misses (V4a, V4b, V5, V5c, V6, V7) are where the papers' stated limits put them; METHOD.md sec. 4 explains each. Fixing them would mean moving past these papers to higher fidelity, which is out of scope by the author's decision of 2026-11-02.
\item PAIN\_Program is not touched. The Goody-in-PAIN idea is parked at the author's request.
\end{itemize}

\section*{Second review (2026-11-02, /code-review)}

{\small
\begin{longtable}{>{\raggedright\arraybackslash}p{0.037\linewidth}>{\raggedright\arraybackslash}p{0.436\linewidth}>{\raggedright\arraybackslash}p{0.436\linewidth}}
\hline
\textbf{\#} & \textbf{Finding} & \textbf{Decision} \\
\hline
\endhead
R1 & The A-weighted view crashed on a static segment, because there is no boundary layer to set the axis base. & \textbf{Fixed.} The base falls back to the total spectrum. The GUI smoke test now runs a static segment. \\
R2 & A zero skin, trim or pane mass made the mass law (eq. 46) 0/0, so NR was NaN. & \textbf{Fixed.} The pipeline raises a clear error. Test: \texttt{test\_\allowbreak{}undefined\_\allowbreak{}inputs\_\allowbreak{}are\_\allowbreak{}rejected}. \\
R3 & Zero cabin absorption made the room equation and \texttt{tau\_\allowbreak{}equivalent} 0/0. & \textbf{Fixed}, by the same check and test. \\
R4 & projectPlan.md had no implementation notes for the audit changes. & \textbf{Fixed.} Notes added. \\
R5 & The Mission and Propeller readouts ignored the measured OAT that the case itself uses. & \textbf{Fixed.} They use the same atmosphere as \texttt{operating\_\allowbreak{}point}. \\
R6 & The case runners built a dummy \texttt{Segment} just to get the atmosphere. & \textbf{Fixed.} They now call \texttt{air(T, isa(alt).\allowbreak{}p)} directly. Results are unchanged. \\
R7 & The validation tab had a dead \texttt{"ka"} axis branch. & \textbf{Removed.} \\
\hline
\end{longtable}}

No result changed: verification is still 4/4, and the validation results are as recorded.

\section*{State after the audit and review}

\begin{itemize}
\item \textbf{Tests:} 48 pass (\texttt{pytest}); \texttt{pyflakes} is clean; the GUI smoke test passes, including the static segment.
\item \textbf{Verification:} 4/4 checks ok.
\item \textbf{Validation:} 1 of 8 met (V5b); V6 at 72 \% (75 \% needed); the others as recorded in cases.yaml.
\item \textbf{Default cruise case:} 98.2 dBA room average, 103.0 dBA at the 95th-percentile seat.
\end{itemize}

\end{document}
